\documentclass[10pt]{article}
\usepackage[margin=20mm]{geometry}
\usepackage{amsmath,amssymb,mathtools,fontspec,unicode-math,xurl,hyperref}
\IfFontExistsTF{DejaVu Serif}{\setmainfont{DejaVu Serif}}{\setmainfont{Latin Modern Roman}}
\setmathfont{latinmodern-math.otf}
\setlength{\parindent}{0pt}\setlength{\parskip}{6pt}\sloppy
\title{A technical review of
Stability of coherent pattern formation through invasion in the FitzHugh–Nagumo system}
\author{AI-assisted technical review}
\date{}
\begin{document}
\maketitle
Abstract: The exact manuscript arXiv:2310.08765v2 is reviewed. Three substantive problems concern the sign of the group velocity in a defining hypothesis, cancellation in the critical resolvent decomposition, and the explicit high-frequency resolvent expansion. The last problem admits a direct counterexample to the stated remainder estimate using compactly supported smooth test functions. An additional regularity estimate is false even on a uniformly bounded sequence of smooth odd inputs. These findings require corrections to hypotheses and core linear arguments. They do not constitute a counterexample to the intended nonlinear stability theorem after correction. The supported assessment is Major for this exact version; scientific adjudication remains pending.\par
\section*{Exact source and review scope}
M25-0297 is by Montie Avery, Paul Carter, Björn de Rijk and Arnd Scheel [paper]. The examined file is arXiv:2310.08765v2, submitted on 24 March 2026 at 19:29:38 UTC, with 83 pages and seven figures. Its SHA256 is `0f68ecc6295f2931c2916281ea410c684ddb33d91ba83a2d61b9df25783900f6`. Physical and printed page numbers agree. EMS associates the title with DOI 10.4171/JEMS/1670, acceptance on 23 March 2025 and online publication on 27 July 2025 [publisher]. Thus this author version postdates the publication date. The journal PDF was unavailable in the acquired set; no equivalence with the journal version is asserted. The frozen 2025 identifier and original title mapping are retained.\par
All 83 pages, including appendices and references, were read. Twenty-nine source pages were actually viewed as images, including all seven figures and the decisive formulas below. Coverage includes the hypotheses, weighted spaces, spectral branches, Fredholm bordering, contour estimates, resolvent decomposition, phase modulation, nonlinear damping, bootstrap and appendices A–E. This is not a formal verification of all external spectral theory or of the companion works.\par
Independent symbolic calculations checked the signs, algebraic decompositions, kernel jumps, polynomial remainders and exact norm formulas used here. Analytical arguments below supply the functional and operator implications. No author code, numerical front solver or Evans-function computation was run. Bounded searches of arXiv, EMS and author publication pages located no separate correction or code release. The related 2026 preprint Diffusive synchronization of phase waves in the FitzHugh–Nagumo system is publicly listed [related, author]; it is not silently identified with the differently titled in preparation reference [4] of the reviewed manuscript, nor was its 103-page proof independently certified. An unlocated supplement or code release is not itself a mathematical defect.\par
\section*{Main statement and dependencies}
The system is \[u_t=u_{xx}+u(u+a)(1-u-a)-w,\qquad
 w_t=\varepsilon(u-\gamma w),\] with \(D=\operatorname{diag}(1,0)\), \(0<a<1/2\), \(0<\gamma<4\), and sufficiently small positive \(\varepsilon\). Theorem 1.1 supplies the pulled front in its specified parameter range. The precise nonlinear result, Theorem 2.2 (pp. 14–15), additionally assumes diffusive spectral stability of the wave train, outgoing group velocity, and absence of unstable point spectrum and a bounded zero resonance (Hypotheses 1–3, pp. 13–14).\par
Its perturbation is small in \(\|\omega w_0\|_{L^1_{0,1}}+\|\omega w_0\|_{H^3\times H^2}\). The conclusions distinguish the unmodulated decay, the faster decay after a spatially varying phase correction, and the improved \(t^{-3/2}\) leading-edge estimate. In particular, stability near an already constructed front is not a proof of selection from all steep or compactly supported initial data. The discussion on p. 56 leaves that selection question for further work.\par
The proof passes from spectral information in Section 3 to the resolvent in Section 6, the linear estimates of Theorem 7.1, and the nonlinear estimates and bootstrap in Sections 8–9. Appendix A supplies the high-frequency contour control used in Proposition 4.3. The issues below are evaluated on these dependencies, rather than by requiring a counterexample to the whole nonlinear conclusion.\par
\section*{Confirmed substantive problems}
Problem 1 (Major: the group-velocity hypothesis has the opposite sign).  Equation (2.14), as invoked explicitly in Hypothesis 2, is incompatible with the Bloch operator and dispersion expansion printed on p. 13.\par
Location and assumptions. The displayed Bloch operator is \[\widehat L_\mathrm{wt}(\nu)
 =D(\partial_\xi+\nu)^2+c_\mathrm{lin}(\partial_\xi+\nu)+F'(u_\mathrm{wt}).\] Under Hypothesis 1, its zero eigenvalue is algebraically simple with right eigenvector \(u_\mathrm{wt}'\). The adjoint vector is normalized by \(\langle u_{\rm ad},u_\mathrm{wt}'\rangle=1\). Equation (2.13) defines \(\lambda_\mathrm{wt}(\nu)=-c_g\nu+D_{\rm eff}^{\mathrm{wt}}\nu^2+O(\nu^3)\); (2.14) instead sets \[c_g=2\langle u_{\rm ad},Du_\mathrm{wt}''\rangle+c_\mathrm{lin}.\] Hypothesis 2 (p. 14) requires the quantity given by (2.14) to be negative.\par
Decisive calculation. Differentiate the simple eigenvalue equation and pair with the adjoint kernel. The differentiated eigenvector term vanishes, yielding \[\lambda_\mathrm{wt}'(0)
 =\langle u_{\rm ad},(2D\partial_\xi+c_\mathrm{lin})u_\mathrm{wt}'\rangle
 =:K=2\langle u_{\rm ad},Du_\mathrm{wt}''\rangle+c_\mathrm{lin}.\] Consequently (2.13) requires \(c_g=-K\), whereas (2.14) gives \(c_g=K\). Together these equalities would force \(K=c_g=0\), excluded by Hypothesis 2. If the hypothesis is read literally as \(K<0\), the spatial root has \(\nu_\mathrm{wt}'(0)=1/K<0\), opposite to the positive coefficient required in Proposition 3.1 (p. 17) and Section 6.1. The direction also follows from the Fourier factor \(e^{ik\xi+\lambda_\mathrm{wt}(ik)t}\): the packet velocity is \(-K\).\par
Impact and repair boundary. This reverses the admissibility condition selecting an outgoing mode; it is not an inconsequential label within the proof. Keeping (2.13) requires replacing (2.14) by \(c_g=-K\), so the intended condition is \(K>0\). Then inversion gives \[\nu_\mathrm{wt}(\lambda)=-\lambda/c_g
                  +D_{\rm eff}^{\mathrm{wt}}\lambda^2/c_g^3+O(\lambda^3),\] consistent with the signs in Proposition 3.1 when \(c_g<0\). The differentiation and inversion have been checked independently. Whether a particular front satisfies the corrected spectral hypotheses must still be established; no front satisfying all the inconsistent printed conventions is claimed as a counterexample.\par
Problem 2 (Major: the critical resolvent residual does not cancel).  The displayed residual in (6.7) fails the identity from which it is defined and therefore lacks the small spectral factor required in Lemma 7.4. The final decomposition (6.28) has a further sign error.\par
Exact algebra and location. On p. 31, let \(K_\lambda g\) denote the integral in (6.5) without its leading minus sign. Let \(K_0(\lambda)g\) denote the same integral with \(Q,Q^{-1},P_\mathrm{wt}^{cu}\) frozen at zero, retaining the exponential \(e^{\nu_\mathrm{wt}(\lambda)(\xi-\zeta)}\). Equation (6.5) gives \[(L_\mathrm{wt}-\lambda)^{-1}g=-K_\lambda g+\bar s_e^\mathrm{wt}(\lambda)g.\] The next display extracts \(-K_0(\lambda)\), so the necessary residual is \(K_0(\lambda)-K_\lambda\). But (6.7) prints \(-K_0(\lambda)-K_\lambda\): the first integrand has a minus sign and the second integral is also subtracted. Reconstruction therefore has error \(-2K_0(\lambda)g\), rather than zero.\par
This is a nonzero operator, not merely a formal possibility. The rank-one neutral spectral projection represents the simple translation mode. Its adjoint pairing with suitable compactly supported forcing is nonzero. Equivalently, the injection \(\Lambda_1 g\) cannot have zero neutral projection for all \(g\): a nonzero adjoint first-order Floquet solution cannot have both its second and third components identically zero, since its adjoint equations would then force its first component to vanish. After output projection, the associated mode is \(u_\mathrm{wt}'\ne0\). Thus the error survives on permitted smooth test data. It does not acquire a factor \(\lambda\).\par
The second integral in Lemma 7.4, p. 45, is estimated by explicitly asserting such an extra factor in (6.7), to gain \(t^{-1/2}\). This assertion fails for the displayed residual. The faster residual decay in Theorem 7.1, p. 41, is then used in (9.9)–(9.13), p. 54.\par
There is a separate exact mismatch in (6.28), p. 39. Write \(r=u_\mathrm{fr}'\), \(w=u_\mathrm{wt}'\), and \(q=q(\cdot,\sigma^2)\). The contribution being split is \(\alpha_-\chi_-e^\nu q\). The principal part is \(\alpha_-\chi_-e^\nu r\); the printed residual uses \((r-w)+(w-q)\). Their sum is \(\alpha_-\chi_-e^\nu(2r-q)\), with error \(2\alpha_-\chi_-e^\nu(r-q)\). The residual must instead use \((w-r)+(q-w)\).\par
Neutral vector and forcing normalization. Lemma 6.10 (p. 37) also prints \(q(\cdot,0)=u_\mathrm{wt}\). The zero mode required by Hypothesis 1 is \(u_\mathrm{wt}'\), as used later on p. 39 and in (7.14), p. 45. For this nonconstant wave, \[L_\mathrm{wt}u_\mathrm{wt}=F'(u_\mathrm{wt})u_\mathrm{wt}-F(u_\mathrm{wt})
       =\big((1-2a)u_{\mathrm{wt},1}^2-2u_{\mathrm{wt},1}^3,\,0\big)^T\ne0.\] Nonvanishing follows because a nonconstant continuous first component cannot take values only in the finite zero set of this polynomial. A constant first component would also force the second component constant through the traveling-wave equations.\par
In the first-order reduction (6.2), p. 30, the stated variables \(U=(u,u_\xi,v)^T\) additionally require \(G=(0,g_1,g_2/c_\mathrm{lin})^T\), because the original second equation has coefficient \(c_\mathrm{lin}\partial_\xi v\). The displayed \(\Lambda_1 g=(0,g_1,g_2)^T\) misses this factor. That local normalization error alone need not change decay exponents, but must be corrected in an exact resolvent identity.\par
Repair and limits. Replacing the residual by \(K_0-K_\lambda\) produces an analytic coefficient difference of order \(\lambda\); reversing both differences in (6.28), using \(q(\cdot,0)=u_\mathrm{wt}'\), and correcting \(\Lambda_1\) restore the corresponding local identities. These algebraic repairs have been checked. A fully propagated proof of all subsequent estimates is not supplied here. Major severity reflects failure of an essential cancellation supporting the main linear and nonlinear estimates, not a demonstrated nonlinear instability.\par
Problem 3 (Major: the explicit high-frequency approximation has the wrong sign).  The kernels in (A.3) and (A.9) represent the negatives of the inverses claimed there. The remainder estimate of Lemma A.4 is false for the explicit terms (A.10)–(A.16) as printed.\par
Kernel identities. Appendix A sets \[\begin{gathered}
 L_0^{11}=\partial_X^2-i\kappa,\qquad
 L_0^{22}=c\partial_X-\widetilde b\mu-i\kappa/\mu+\mu B_0,\\
 c=c_\mathrm{lin}>0,\qquad B_0\le0,\qquad\widetilde b=b+\varepsilon\gamma>0.
 \end{gathered}\] For \(k=\sqrt{i\kappa}\) with positive real part, (A.9), p. 60, uses \(G^p(X)=e^{-k|X|}/(2k)\). Its derivative jump at zero is \(-1\). It solves \((\partial_X^2-i\kappa)G^p=-\delta_0\), not \(+\delta_0\). Likewise the kernel (A.3), p. 58, including its oscillatory factor, equals \(1/c\) at \(X=Y\) from the left and zero from the right. The coefficient of the delta distribution after applying \(c\partial_X\) is \(-1\). Away from the diagonal the homogeneous equation holds, including the actual \(B_0\) term. Thus it too is the negative inverse. These are identities for the printed operators and allowed coefficients, not counterexamples using a different model.\par
Counterexample to Lemma A.4. The following uses the actual smooth bounded front coefficients. Fix \(b>0\) within the permitted strip and \(0\ne\phi\in C_c^\infty(\mathbb R,\mathbb R)\). Let \(\lambda=b+i\Omega\), \(\Omega\to+\infty\), and define \[U=(0,\phi)^T,\qquad
 g_\Omega=(L_\mathrm{fr}-\lambda)U
   =\bigl(-\phi,\,
       c\phi'+(b_0-\varepsilon\gamma-b-i\Omega)\phi\bigr)^T.\] These data are compactly supported and smooth, and \(\|g_\Omega\|_{L^p}=O(\Omega)\), for example with \(p=2\). For sufficiently large \(\Omega\) the resolvent exists: the diagonal inversion and small block-coupling argument of Appendix A applies with the corrected inverse signs; those signs do not affect the norm bounds. Hence \((L_\mathrm{fr}-\lambda)^{-1}g_\Omega=U\).\par
Let \(T_\Omega\) be the positive transport kernel printed in (A.11)–(A.12), and \(H_\Omega\) the positive heat kernel in (A.10). The jump calculation gives exactly \[T_\Omega\bigl(c\partial_\xi+b_0-\varepsilon\gamma-\lambda\bigr)\phi
 =-\phi.\] The exponential bound \(b_0\le0\) implies \(\|T_\Omega\|_{L^p\to L^p}\le(b+\varepsilon\gamma)^{-1}\); the heat-kernel \(L^1\) bound implies \(\|H_\Omega\|_{L^p\to L^p}\le C/\Omega\). Consequently the printed formulas give \[\begin{aligned}
 v_0&=T_\Omega g_{\Omega,2}=-\phi,\\
 v_1&=-\varepsilon T_\Omega H_\Omega g_{\Omega,1}=O_{L^p}(\Omega^{-1}),\\
 v_2&=-\varepsilon T_\Omega H_\Omega T_\Omega g_{\Omega,2}
       =O_{L^p}(\Omega^{-1}).
 \end{aligned}\] The second component of the remainder in Lemma A.4 therefore tends to \(2\phi\) in \(L^p\), while its asserted upper bound is \(C\Omega^{-3/2}\|g_\Omega\|_{L^p}=O(\Omega^{-1/2})\to0\). This is a contradiction. If the lemma’s real-valued test-space notation is enforced, its claimed bound on all real data implies the complex-data bound by splitting real and imaginary parts, with at most a fixed factor in the constant. The same contradiction applies.\par
Dependence and repair boundary. Lemma A.4 is used for the contour shift and Proposition A.5, hence Proposition 4.3 and the high-frequency part of Theorem 7.1. Each diagonal Green kernel needs a leading minus sign. Terms with two such inverses, such as \(u_1,v_1\), keep their sign; terms with an odd number, including \(u_0,v_0,v_2\), reverse sign. A global negation of every displayed term is not the repair. The corrected Neumann-series normalization is explicit, but the complete contour proof has not been recertified here. The actual false remainder estimate, rather than the ease of changing a symbol, determines Major severity.\par
\section*{Additional confirmed regularity issue}
Problem 4 (Minor: the \(L^1\)-to-\(W^{2,\infty}\times W^{1,\infty}\) estimate is false).  Lemma 6.7, (6.17), p. 34, cannot hold for every odd \(g\in L^1_{0,1}\) as stated.\par
Fix any sufficiently small positive real \(\sigma\) in its stated sector. The constant-coefficient differential operator \(L_+-\sigma^2\) maps \(W^{2,\infty}\times W^{1,\infty}\) boundedly into \(L^\infty\). Take \(g=(h,0)^T\), where \[h(\xi)=\operatorname{sgn}(\xi)|\xi|^{-1/2}
                 {\bf1}_{0<|\xi|<1}.\] This is odd, lies in the required weighted \(L^1\) space, and does not belong to \(L^\infty\). A solution in the claimed target space would make \(g=(L_+-\sigma^2)u\) bounded, a contradiction in distributions. For reference, \(\int|h|=4\) and \(\int|\xi h|=4/3\).\par
Restricting to smooth inputs does not restore a uniform estimate. Set \(h_n(\xi)=n^2\xi e^{-n^2\xi^2}\). These are odd Schwartz functions, with \[\|h_n\|_1=1,\qquad
 \int|\xi h_n(\xi)|\,d\xi=\frac{\sqrt\pi}{2n},\qquad
 \|h_n\|_\infty=\frac{n}{\sqrt{2e}}.\] Thus their \(L^1_{0,1}\) norms remain bounded, while composition of the claimed resolvent bound with \(L_+-\sigma^2\) would give a uniform \(L^\infty\) bound on \(h_n\). The contradiction uses all the printed input hypotheses, including oddness.\par
In the first-order Green representation, differentiating once introduces the delta term in \(\partial_\xi T=MT+\delta_0 I\). The highest derivatives therefore retain the unsmoothed forcing. The discussion around (B.3), pp. 69–70, cannot justify the stated extra derivatives by ignoring that jump.\par
The natural weaker target is \(W^{1,\infty}\times L^\infty\); highest derivatives require an additional \(L^\infty\) norm of the input. The commutator in Lemma 6.9 needs only the first derivative of the first component, and Lemma 6.16 needs only the \(L^\infty\) bound; the separate \(L^2\) statement is Lemma 6.8. Accordingly this overstatement does not, by itself, establish failure of the main decay conclusion. Minor severity is based on that restricted downstream use. The local kernel argument supports the weaker target in the prescribed sector; no estimate in a larger spectral region is claimed.\par
\section*{Other checked points and limits}
Problem 5 (Minor: local weight sign and exponential norm power).  Lemma 3.2 (p. 18) has the wrong sign for the leading-edge weight, and the heuristic following (5.1), p. 23, gives the wrong power for the half-line \(L^2\) norm.\par
Lemma 3.2 (p. 18) prints \(\eta_\mathrm{lin}<0\), whereas the growing weight and the tail \(e^{-\eta_\mathrm{lin}\xi}\) require the positive branch. At \(\varepsilon=0\) the double root gives \(\eta_\mathrm{lin}=\sqrt{a(1-a)}>0\); the branch remains positive for sufficiently small positive \(\varepsilon\). Also the heuristic \(L^2\) blowup following (5.1), p. 23, uses \(|\sigma|^{-1}\) for the norm. Directly, \(\int_0^\infty e^{-2\sigma\xi}\,d\xi=(2\sigma)^{-1}\), so the norm scales as \(|\sigma|^{-1/2}=|\lambda|^{-1/4}\) in a fixed sector. The later Lemmas 6.8 and 6.15 use the correct half-power. These local corrections are recorded as Minor; no extra core failure is inferred from them.\par
The inverse modulation algebra, the quadratic/cubic FitzHugh–Nagumo remainder, and the role of the dissipative second component were checked against the displayed equations. In particular, \[f(u+z)-f(u)-f'(u)z=(1-2a-3u)z^2-z^3.\] The phase-coordinate construction and nonlinear damping are read with their small-derivative and regularity assumptions; they are not declared valid merely because their formal exponents close. The convolution powers in (9.9)–(9.13) rely on the corrected faster linear residual bounds, precisely where the above issues matter.\par
The review establishes the local identities, the explicit failure of the printed Appendix A remainder, and the regularity counterexamples. It has not computed an unstable front, proved failure of the intended nonlinear stability theorem, independently proved Hypotheses 1–3 for the full parameter range, or propagated every proposed correction through all 83 pages. No numerical evidence is substituted for these missing tasks. The exact manuscript requires substantive revision; scientific adjudication remains pending and human reviews are empty. Current independent acceptance is recorded in the associated provenance.\par
9 M. Avery, P. Carter, B. de Rijk and A. Scheel, Stability of coherent pattern formation through invasion in the FitzHugh–Nagumo system, arXiv:2310.08765v2, 24 March 2026, https://arxiv.org/abs/2310.08765v2. EMS Press, article record, DOI 10.4171/JEMS/1670, https://ems.press/journals/jems/articles/14298982 (accessed 27 October 2026). M. Avery, P. Carter, B. de Rijk and A. Scheel, Diffusive synchronization of phase waves in the FitzHugh–Nagumo system, arXiv:2601.05377v1, 8 January 2026, https://arxiv.org/abs/2601.05377v1 (metadata only; full proof not reviewed here). P. Carter, publications page, https://webapps.math.uci.edu/\textasciitilde{}pacarter/papers.html (accessed 27 October 2026).\par
\end{document}
